% !TeX root = ../complete.tex
\section{مسیرهای جریان و برآورد \lr{Static Power}}
\label{sec:power-static}

\subsection{مصرف توان در یک وضعیت منطقی ثابت}
\label{sec:power-static-sources}

در مدل \lr{Switch} یک گیت مکمل \lr{CMOS}، پس از پایان گذار، یکی
از شبکه‌های \lr{Pull-up} یا \lr{Pull-down} خاموش است و مسیر مستقیم
روشنی میان تغذیه و زمین وجود ندارد. خازن خروجی نیز پس از رسیدن به
ولتاژ نهایی جریان شارژ نمی‌گیرد. در ترانزیستور واقعی، حتی در این وضعیت
جریان‌های کوچک از چند مسیر عبور می‌کنند. برای تغذیه ثابت، اگر این
جریان‌ها سهم‌های جداگانه‌ی جریان کشیده‌شده از همان تغذیه باشند، برآورد
توان ساکن از رابطه زیر به دست می‌آید
\cite[رابطه \secnum{۵.۱۲} و بخش \secnum{۵.۳.۱}،
صفحات ۱۸۵ و ۱۹۴ تا ۱۹۷]{weste}\cite[اسلایدهای ۱۰ و ۲۰]{powerslides}.
\begin{equation}
  P_{\mathrm{stat}}
    =V_{DD}\bigl(I_{\mathrm{sub}}+I_{\mathrm{gate}}+
                 I_{\mathrm{junct}}+I_{\mathrm{cont}}
           \bigr).
  \label{eq:power-static}
\end{equation}
هر جمله داخل پرانتز جریان است و ضرب جمع جریان‌ها در ولتاژ، توان
را برحسب وات می‌دهد. مسیرهای فیزیکی نشتی در
\NotesTopic{sec:leakage}{فصل مدل ترانزیستور} بررسی شده‌اند. در اینجا
اثر هر مسیر را بر جریان کشیده‌شده از منبع دنبال می‌کنیم.

\begin{itemize}
  \item \(I_{\mathrm{sub}}\) سهم \lr{Subthreshold Leakage} است.
    زیر آستانه، ترانزیستور خاموش هنوز می‌تواند جریان کوچکی از مسیر
    کانال عبور دهد. ولتاژهای پایانه‌ها و دما بر این جریان اثر دارند.
    در چند ترانزیستور خاموشِ سری، ولتاژ گره میانی از تعادل جریان
    به دست می‌آید و بایاس هر ترانزیستور با حالت تک‌ترانزیستور فرق
    می‌کند. کاهش نشتی حاصل در بسیاری از وضعیت‌های ورودی را
    \lr{Stack Effect} می‌نامیم.
  \item \(I_{\mathrm{gate}}\) سهم \lr{Gate Leakage} است. اختلاف
    پتانسیل در دو سوی دی‌الکتریک نازک گیت می‌تواند موجب عبور حامل‌ها
    به‌واسطه‌ی \lr{Tunneling} شود. ضخامت و جنس دی‌الکتریک و ولتاژ
    دو سوی آن تعیین می‌کنند این سهم چقدر باشد.
  \item \(I_{\mathrm{junct}}\) سهم \lr{Junction Leakage} است.
    ناحیه‌های \lr{Source} و \lr{Drain} با \lr{Body} پیوند تشکیل
    می‌دهند. حتی یک پیوند با بایاس معکوس جریان کوچکی دارد. میدان
    زیاد نزدیک پیوند می‌تواند سهم سازوکارهایی مانند
    \lr{Band-to-Band Tunneling} را بیشتر کند.
  \item \(I_{\mathrm{cont}}\) سهم \lr{Contention Current} است.
    برای نمونه، در \lr{Pseudo-nMOS} یک \lr{pMOS} همیشه روشن
    خروجی را به سمت تغذیه می‌کشد. هنگامی که شبکه \lr{nMOS} خروجی
    را پایین نگه می‌دارد، هر دو مسیر هم‌زمان هدایت می‌کنند و جریان
    ماندگار برقرار است. این جریان در مدل ماندگار گیت مکمل
    \lr{CMOS} وجود ندارد. \lr{Short-Circuit Current} بخش
    \ref{sec:power-short-circuit} پالس جریان در زمان گذار است.
\end{itemize}

\subsection{برآورد نشتی از عرض افزاره‌ها}
\label{sec:power-static-estimate}

وقتی مقدار یک نشتی به‌ازای واحد عرض معلوم باشد، نخست ترانزیستورها
را برحسب نوع و وضعیت بایاس گروه‌بندی می‌کنیم. در هر گروه، تعداد
افزاره‌ها ضرب‌در عرض متوسط، عرض کل گروه را می‌دهد. اگر کسری از
عرض کل در وضعیت دارای نشتی موردنظر باشد، همان کسر را در عرض کل
و جریان بر واحد عرض ضرب می‌کنیم. برای مثال، در مدل ساده‌ی نشتی
زیرآستانه، \(W_g\) عرض کل گروه \(g\)، \(b_g\) کسر عرضِ خاموش
و \(j_{\mathrm{off},g}\) جریان خاموش بر واحد عرض آن گروه است.
\cite[بخش \secnum{۵.۳.۱.۵}، صفحه ۱۹۷]{weste}
\begin{equation}
  I_{\mathrm{sub}}\approx\sum_g b_g W_g j_{\mathrm{off},g}.
  \label{eq:power-width-leakage}
\end{equation}
\(b_g\) بی‌بعد است. اگر \(W_g\) برحسب میکرومتر و
\(j_{\mathrm{off},g}\) برحسب نانوآمپر بر میکرومتر باشد، حاصل ضرب
برحسب نانوآمپر خواهد بود. ثبات مقدار جریان بر واحد عرض در هر گروه
یک تقریب است. به‌ویژه، جدول \secnum{۵.۲} کتاب نشان می‌دهد که در
یک \lr{NAND} سه‌ورودی، تغییر ترتیب ورودی‌ها می‌تواند ولتاژ گره‌های
میانی و درنتیجه نشتی را عوض کند \cite[صفحه ۱۹۶]{weste}.

در مثال ۵.۴ کتاب و اسلایدهای ۲۱ و ۲۲، دوباره همان تراشه‌ی بخش
\ref{sec:power-effective-capacitance} را بررسی می‌کنیم. از ۵۰ میلیون
ترانزیستور منطق با عرض متوسط \qty{0.30}{\micro\metre}، سهم
۵ درصد در گروه با آستانه پایین‌تر و ۹۵ درصد در گروه آستانه بالا
قرار می‌گیرند. اسلاید گروه ۵ درصدی را \lr{Normal \(V_t\)} و
کتاب در مثال ۵.۴ آن را \lr{low-threshold} می‌نامد. هر دو منبع
برای این گروه در حالت خاموش، مقدار
\qty[per-mode=symbol]{100}{\nano\ampere\per\micro\metre}
به کار می‌برند. \(W_{\mathrm{nom}}\) عرض کل این گروه و
\(W_{\mathrm{high}}\) عرض کل گروه آستانه بالاست.
همه‌ی ۹۵۰ میلیون ترانزیستور حافظه با عرض متوسط
\qty{0.10}{\micro\metre} در گروه آستانه بالا قرار می‌گیرند.
عرض کل دو گروه از جمع عرض ترانزیستورهایشان به دست می‌آید
\cite[مثال ۵.۴، صفحه ۱۹۷]{weste}\cite[اسلایدهای ۲۱ و ۲۲]{powerslides}.
\[
  \begin{aligned}
    W_{\mathrm{nom}}&=(50\times10^6)(0.05)
                       (\qty{0.30}{\micro\metre})
                     =\qty{0.75e6}{\micro\metre},\\
    W_{\mathrm{high}}&=(50\times10^6)(0.95)
                        (\qty{0.30}{\micro\metre})
                        +(950\times10^6)(\qty{0.10}{\micro\metre})\\
                    &=\qty{109.25e6}{\micro\metre}.
  \end{aligned}
\]
در این برآورد، نصف عرض هر گروه در حالت خاموش است. \lr{Stack Effect}
را کنار می‌گذاریم و با مقادیر خاموشِ داده‌شده در مثال، سهم کانال را
حساب می‌کنیم. برای سهم گیت فرض می‌کنیم نیمه‌ی روشنِ عرض کل دارای
آن نشتی است. پس نصف مجموع عرض را در
\qty[per-mode=symbol]{5}{\nano\ampere\per\micro\metre} ضرب می‌کنیم.
\[
  \begin{aligned}
    I_{\mathrm{sub}} &\approx \tfrac12\bigl[
       W_{\mathrm{nom}}(\qty[per-mode=symbol]{100}{\nano\ampere\per\micro\metre})
      +W_{\mathrm{high}}(\qty[per-mode=symbol]{10}{\nano\ampere\per\micro\metre})
      \bigr]\approx\qty{584}{\milli\ampere},\\
    I_{\mathrm{gate}}&\approx\tfrac12(W_{\mathrm{nom}}+W_{\mathrm{high}})
        (\qty[per-mode=symbol]{5}{\nano\ampere\per\micro\metre})
        =\qty{275}{\milli\ampere}.
  \end{aligned}
\]
با فرض ناچیز بودن نشتی پیوند و نبود مسیر ماندگار ناشی از هدایت
هم‌زمان، جمع این دو جریان در تغذیه \qty{1.0}{\volt} حدود
\qty{859}{\milli\watt} توان ساکن می‌دهد. این مقدار در همان مدل
تراشه حتی اگر ساعت متوقف شود و وضعیت بایاس و دما ثابت بمانند،
با صفر شدن توان سوئیچینگ از بین نمی‌رود. نسبت آن به برآورد
\qty{6.1}{\watt} سوئیچینگ بخش \ref{sec:power-effective-capacitance}
حدود ۱۴ درصد است. هر دو برآورد به داده‌ها و فرض‌های همین مثال
وابسته‌اند.

\subsection{کاهش نشتی با انتخاب افزاره و وضعیت ورودی}
\label{sec:power-leakage-controls}

مقدار \(V_t\) تعیین می‌کند ترانزیستور در یک ولتاژ گیت معین چه
مقدار جریان عبور می‌دهد. با ثابت بودن سایر بایاس‌ها و دما، بالا بردن
آستانه به اندازه \(\Delta V_t\) جمله نماییِ جریان زیرآستانه را
در مدل فصل ۲ در عامل زیر ضرب می‌کند. \(S\) تغییر ولتاژ لازم برای
یک دهه تغییر جریان است
\cite[بخش‌های \secnum{۲.۴.۴} و \secnum{۵.۳.۱.۱}]{weste}.
\[
  \frac{I_{\mathrm{sub,new}}}{I_{\mathrm{sub,old}}}
  \approx 10^{-\Delta V_t/S}.
\]
برای مثال، اگر \(S\) برابر \qty{100}{\milli\volt} بر دهه باشد،
افزایش \qty{50}{\milli\volt} آستانه، نمای کسر را به \(-1/2\)
می‌رساند. \(10^{-1/2}=1/\sqrt{10}\) است، پس جریان تخمینی به حدود
۳۲ درصد مقدار اولیه می‌رسد. مقدار \(S\) و جریان اولیه از بایاس، دما
و فناوری می‌آیند.

در کتابخانه دارای چند آستانه، \lr{High-\(V_t\)} برای مسیر دارای
مهلت کافی، نشتی زیرآستانه‌ی کمتری دارد. \lr{Low-\(V_t\)} در بخش
حساس به تأخیر می‌تواند جریان روشن بیشتری بدهد. دی‌الکتریک ضخیم‌تر
در فناوری‌هایی که آن گزینه را دارند، می‌تواند نشتی گیت را کم کند.
سرعت و سازگاری سلول با ولتاژ کاری هنگام انتخاب هر نوع بررسی می‌شوند
\cite[بخش \secnum{۵.۳.۳}، صفحه ۱۹۹]{weste}.

در \lr{Body Bias}، ولتاژ \lr{Body} نسبت به \lr{Source} تنظیم
می‌شود. در \lr{nMOS}، افزایش بایاس معکوس پیوند بدنه و سورس در
محدوده معتبر مدل، آستانه را بالا می‌برد و می‌تواند جریان زیرآستانه
را کم کند. بایاس معکوس زیاد، میدان پیوند را افزایش می‌دهد و ممکن
است نشتی پیوند بیشتر شود. با \lr{Input Vector Control} نیز می‌توان
هنگام بی‌کاری ترکیبی از ورودی‌ها را اعمال کرد که وضعیت ترانزیستورهای
خاموش در شاخه‌های سری نشتی کمتری ایجاد کند. هزینه تغییر ورودی‌ها
و صرفه‌جویی در طول بی‌کاری هر دو در ارزیابی این روش سهم دارند
\cite[بخش‌های \secnum{۵.۳.۴} و \secnum{۵.۳.۵}،
صفحات ۱۹۹ و ۲۰۰]{weste}.
